Se estudia el problema del robot explorador: dada una grilla de $N\times M$ celdas con muros, minas y celdas con puntaje, el robot parte en la esquina superior izquierda con $V$ vidas y debe llegar a la esquina inferior derecha moviéndose solo a la derecha o hacia abajo, maximizando el puntaje; pisar una mina cuesta una vida y una penalización $P$. Se comparan fuerza bruta con poda, programación dinámica \textit{bottom-up} \cite{cormen} y una heurística greedy \cite{algorithms_erickson}, midiendo tiempo, memoria y calidad de la solución.
